% !TeX program = pdflatex
\documentclass[11pt,a4paper]{article}
\usepackage[margin=1in]{geometry}
\usepackage{amsmath,amsthm,mathtools}

\usepackage[T1]{fontenc}
\usepackage[utf8]{inputenc}
\usepackage{newtxtext,newtxmath}
\ifdefined\pdfmapfile\pdfmapfile{+newtx.map}\fi
\usepackage[expansion=false]{microtype}
\usepackage{booktabs}
\usepackage{xcolor}
\definecolor{linkblue}{RGB}{36,73,112}
\usepackage[colorlinks=true,linkcolor=linkblue,citecolor=linkblue,urlcolor=linkblue]{hyperref}
\newcommand{\Eqref}[1]{\hyperref[#1]{Eq.~(\ref*{#1})}}
\newcommand{\Secref}[1]{\hyperref[#1]{Section~\ref*{#1}}}
\newcommand{\Tabref}[1]{\hyperref[#1]{Table~\ref*{#1}}}
\newcommand{\TV}{\mathrm{TV}}
\theoremstyle{plain}
\newtheorem{proposition}{Proposition}[section]
\hypersetup{pdftitle={Research Directions for Variable Selection in Wasserstein Robust Testing},
 pdfauthor={Research note}}
\setlength{\parindent}{1.3em}
\setlength{\parskip}{0pt}
\setlength{\emergencystretch}{2em}
\clubpenalty=10000
\widowpenalty=10000
\setlength{\jot}{2pt}
\AtBeginDocument{
\setlength{\abovedisplayskip}{6pt plus 2pt minus 1pt}
\setlength{\belowdisplayskip}{6pt plus 2pt minus 1pt}}
\renewenvironment{proof}
{\par\pushQED{\qed}\normalfont\noindent\textit{Proof.}\ \ignorespaces}
{\popQED\par}
\title{Research Directions for Variable Selection in Wasserstein Robust Testing}
\author{}
\date{}
\makeatletter
\renewcommand{\maketitle}{
\begin{center}
{\Large\bfseries\@title\par}
\end{center}\vspace{2pt}}
\makeatother

\begin{document}
\maketitle

\section{Starting point}\label{sec:start}
The theory note writes the selection problem
$\min_{|S|\le d}V(S)$ as the minimization of a convex function over
$\mathcal Z_d$ and solves it by outer approximation. In development
checks with the hinge loss, the method evaluated every subset before
it closed the gap, for every cut we tried, and the relaxation was
weak: at the uniform point $z=(d/D)\mathbf1$ it took the value $0.58$,
against the optimum $1.40$ for $D=10$, $d=2$.
This note argues that the weakness belongs to the problem, not to the
formulation, and builds the research directions on a reformulation
that makes this visible.

The argument runs in one line. \Secref{sec:divergence} rewrites
$V(S)$ as a divergence between the projected classes; for the hinge
loss it becomes a budgeted matching. \Secref{sec:gap} uses the
matching to explain the relaxation gap and the hardness.
\Secref{sec:directions} turns the three consequences into a
statistical, a computational and a game-theoretic direction, and
\Secref{sec:priorities} ranks them.

Notation follows the theory note. For a distribution $Q$ on
$\mathbb R^D$, $Q[S]$ is its image under $x\mapsto x[S]$, and
$\mathcal P_k[S]=\{Q:\mathcal W_1(\widehat{\mathbb P}_k[S],Q)\le\rho_k\}$
is the projected ambiguity set, with the $\ell_1$ cost on
$\mathbb R^{|S|}$.

\section{The robust value is a divergence between projections}
\label{sec:divergence}
For $a,b\ge0$ let $\psi(a,b)=\inf_{t\in\mathbb R}\{a\ell(-t)+b\ell(t)\}$.
It is concave and positively homogeneous, and $\psi(a,a)=2a$ because
convexity gives $\ell(-t)+\ell(t)\ge2\ell(0)=2$. Hence
\[
D_\ell(Q_0,Q_1)=2-\int\psi(dQ_0,dQ_1)
\]
is an $f$-divergence between probability distributions~\cite{nguyen2009surrogate}.
For $\ell(t)=(1+t)_+$, $e^t$ and $\log_2(1+e^t)$ it equals $2\,\TV$,
the squared Hellinger distance $\int(\sqrt{dQ_0}-\sqrt{dQ_1})^2$, and
twice the Jensen--Shannon divergence in bits.

\begin{proposition}\label{prop:divergence}
For every $S\subseteq[D]$,
\[
V(S)=2-\min_{Q_k\in\mathcal P_k[S]}D_\ell(Q_0,Q_1).
\]
\end{proposition}
\begin{proof}
Fix $Q_0$ and $Q_1$. Minimizing the risk
$\mathbb E_{Q_0}\ell(-f)+\mathbb E_{Q_1}\ell(f)$ point by point over
$f\in\mathcal F_S$ gives $\int\psi(dQ_0[S],dQ_1[S])$. Projection maps the
ambiguity sets into $\mathcal P_k[S]$, and every element of
$\mathcal P_k[S]$ lifts with the same cost. Exchanging the minimum over
$f$ and the supremum over distributions therefore gives
$V(S)\ge2-\inf D_\ell$, and Lemma~2.1 of the theory note attains this
bound with distributions supported on the projected samples.
\end{proof}

Variable selection is thus the choice of $d$ coordinates on which the
two classes stay as far apart as possible, in the divergence chosen by
the loss, after an adversary moves each class by $\rho_k$ in
$\mathcal W_1$. This places the problem next to methods that select a
projection maximizing a distance between empirical distributions:
sparse principal differences analysis~\cite{mueller2015principal},
projected Wasserstein tests~\cite{wang2021projected}, feature-robust
transport~\cite{petrovich2022feature} and kernel variable
selection~\cite{wang2023variable}. Two things differ here: the
distance is a divergence measured after a transport attack, and its
value is the risk certificate of a minimax test. The mirror
construction, a Wasserstein distance measured after a total-variation
attack, is the robust transport of~\cite{nietert2023robust}.

For the hinge loss the divergence is total variation, and the attack
becomes combinatorial. Write
$c_{ij}(S)=\|x_i^{(0)}[S]-x_j^{(1)}[S]\|_1$ and
\begin{equation}\label{eq:matching}
m_S(\rho)=\max_{\pi\ge0}\Bigl\{\sum_{i,j}\pi_{ij}:\
\sum_j\pi_{ij}\le\frac1{n_0},\ \sum_i\pi_{ij}\le\frac1{n_1},\
\sum_{i,j}\pi_{ij}c_{ij}(S)\le\rho\Bigr\},
\end{equation}
the largest cross-class mass that an adversary can pair up with total
transport budget $\rho$.

\begin{proposition}[Hinge loss]\label{prop:matching}
$V(S)=2\,m_S(\rho_0+\rho_1)$.
\end{proposition}
\begin{proof}
By \hyperref[prop:divergence]{Proposition~\ref*{prop:divergence}},
$V(S)/2=1-\min\TV(Q_0,Q_1)$ is the largest overlap
$\int\min(dQ_0,dQ_1)$. Overlap at a point draws class-$0$ mass from
some $x_i^{(0)}[S]$ and class-$1$ mass from some $x_j^{(1)}[S]$.
Pairing the two, the triangle inequality charges at least $c_{ij}(S)$
per unit of overlap to the two budgets together, so the overlap is at
most $m_S(\rho_0+\rho_1)$. Conversely, let $\pi$ be feasible in
\Eqref{eq:matching} with $\rho=\rho_0+\rho_1$. Move the fraction
$\theta=\rho_0/(\rho_0+\rho_1)$ of each $\pi_{ij}$ from $x_i^{(0)}$ to
$x_j^{(1)}$, and the remaining fraction of it from $x_j^{(1)}$ to
$x_i^{(0)}$. Class $0$ spends at most $\theta\rho=\rho_0$, class $1$
at most $\rho_1$, and the overlap is at least $\sum_{i,j}\pi_{ij}$.
\end{proof}

Three consequences follow, and the rest of the note uses them.
First, the radii enter only through $\rho_0+\rho_1$; on random
instances with unequal radii the solver of the theory note agrees with
$2m_S$ to $10^{-10}$. Second, the least favorable distributions pair up
cross-class points: the adversary is a matchmaker with a transport
budget, and $S$ separates the classes when few cross-class pairs are
cheap in the coordinates of $S$. Third, evaluating $V(S)$ is a
bipartite matching with one side constraint and $n_0n_1$ variables,
instead of a transport plan over $N^2$ pairs with two budgets.

\section{Why convex relaxations are weak}\label{sec:gap}
Let $C^m_{ij}=|x_i^{(0)}[m]-x_j^{(1)}[m]|$, so that
$c_{ij}(S)=\sum_{m\in S}C^m_{ij}$. Linear programming duality for
\Eqref{eq:matching} gives
\[
m_S(\rho)=\min_{u,v,\lambda\ge0}\Bigl\{\rho\lambda+\bar u+\bar v:\
u_i+v_j+\lambda c_{ij}(S)\ge1\Bigr\},\qquad
\bar u=\frac1{n_0}\sum_iu_i,\ \ \bar v=\frac1{n_1}\sum_jv_j.
\]
Each cross-class pair is either paid for by its endpoints or separated
by its priced distance $\lambda c_{ij}(S)$, and selecting $S$ decides
which coordinates the adversary pays for. With $w_m=\lambda z_m$ the
whole selection problem becomes one compact mixed-integer program,
with $\rho=\rho_0+\rho_1$:
\begin{equation}\label{eq:milp}
\begin{aligned}
\min_{z\in\mathcal Z_d,\,u,v,w,\lambda\ge0}\quad&
\rho\lambda+\bar u+\bar v\\
\text{subject to}\quad&
u_i+v_j+\sum_m w_mC^m_{ij}\ge1,\qquad
w_m\le\lambda,\qquad \rho w_m\le z_m,\qquad
\sum_m w_m\le d\lambda.
\end{aligned}
\end{equation}
Its optimal value is $\min_{|S|\le d}m_S(\rho)=\frac12\min_{|S|\le d}V(S)$: the point
$u=\mathbf1$, $v=0$, $\lambda=0$ has value $1$, so $\rho\lambda\le1$
at the optimum, and then $w_m=\lambda z_m$ is optimal.

The relaxation has a plain meaning. Replace $\mathcal Z_d$ by its
convex hull $\{z\in[0,1]^D:\mathbf1^\top z\le d\}$ and set
$z=w/\lambda$. The relaxed value is the smallest matching value over
weighted costs $\sum_mz_m|x[m]-y[m]|$: the relaxation lets the
defender choose a weighted $\ell_1$ metric instead of a set of
coordinates. Weighted metrics are genuinely better defenses. Averaging
coordinates concentrates the pairwise costs $\sum_mz_mC^m_{ij}$ around
their mean, and the adversary loses the few cheap pairs it can always
find in a sparse projection. \Tabref{tab:gap} shows the effect: the
relaxed value is one half to two thirds of the optimum. It is still
stronger than the convex function of the theory note, which already
takes the value $0.58$ at the uniform point on the first instance. The gap
therefore belongs to the problem: the relaxed problem is itself a
sensible robust test, and it is better than every sparse one.

\begin{table}[htbp]
\centering\small
\begin{tabular}{rrrcccc}
\toprule
$D$ & $N$ & $d$ & $\min_SV(S)$ & relaxation of \Eqref{eq:milp} &
\Eqref{eq:milp} by HiGHS & enumeration\\
\midrule
10 & 30 & 2 & 1.40 & 0.89 & 0.2\,s, 17 nodes & 0.2\,s\\
15 & 30 & 2 & 1.12 & 0.75 & 0.4\,s, 25 nodes & 0.5\,s\\
12 & 30 & 3 & 0.87 & 0.49 & 0.7\,s, 302 nodes & 1.1\,s\\
30 & 40 & 3 & 0.75 & 0.43 & 26\,s, 4.1k nodes & 22\,s\\
50 & 40 & 3 & 0.80 & 0.40 & 235\,s, 26k nodes & --\\
100 & 40 & 3 & $\le0.88$ & 0.37 & time limit 600\,s & --\\
\bottomrule
\end{tabular}
\caption{Hinge loss, $\rho_0+\rho_1=0.6$. Both classes are standard
normal except that class~$1$ is shifted by $0.8$ in coordinates
$1$--$3$ and scaled by $1.8$ in coordinate~$4$. Program
\Eqref{eq:milp} reproduces the enumerated optimum wherever enumeration
was run. Development check, not a formal result.}
\label{tab:gap}
\end{table}

The combinatorial core is also hard in the worst case.

\begin{proposition}\label{prop:hard}
Let the data be binary and $\rho_0+\rho_1<\min\{1/n_0,1/n_1\}$. Then
$\min_{|S|\le d}V(S)\le2(\rho_0+\rho_1)$ if and only if some $S$ with
$|S|\le d$ separates every cross-class pair, that is,
$x_i^{(0)}[S]\ne x_j^{(1)}[S]$ for all $i,j$. Deciding this is the
$k$-feature-set problem, which is NP-complete~\cite{davies1994np} and
W[2]-complete in $d$~\cite{cotta2003kfeature}. Hence computing
$\min_{|S|\le d}V(S)$ is NP-hard.
\end{proposition}
\begin{proof}
Write $\rho=\rho_0+\rho_1$. If $S$ separates every pair, then
$c_{ij}(S)\ge1$ for all pairs and $m_S(\rho)\le\rho$. Otherwise some
pair has $c_{ij}(S)=0$ and can be matched at no cost, so
$m_S(\rho)\ge\min\{1/n_0,1/n_1\}>\rho$.
Proposition~\ref{prop:matching} completes the proof.
\end{proof}

Diminishing returns do not rescue greedy selection either. With class
$0$ at $(0,0),(1,1)$ and class $1$ at $(0,1),(1,0)$, each coordinate
alone leaves the projected classes identical, so $V=2$, while the two
coordinates together give $V=2(\rho_0+\rho_1)$ when
$\rho_0+\rho_1\le1$. A strong convex
relaxation is therefore not to be expected. Progress must come from
information that the weighted-metric problem does not see: the sample
size, the matching structure, or the value of mixing coordinates.

\section{Directions}\label{sec:directions}

\subsection{Selection makes robust tests honest in high dimension}
\label{sec:honest}
\Tabref{tab:noise} draws both classes from the same distribution, so
the correct robust value is $2$. Yet $V([D])$ falls to $0.03$ at
$D=50$. In high dimension the empirical points are isolated, every
cross-class pair is expensive to match, and the robust test certifies
a separation that does not exist. Selection limits the effect, but the
best of many subsets is optimistic too, and in the first instance of
\Tabref{tab:gap} the optimal pair consists of two uninformative
coordinates. A radius that ignores the dimension is not honest; an
honest radius must depend on $d$ and on the number of candidate
subsets.

\begin{table}[htbp]
\centering\small
\begin{tabular}{lccccccc}
\toprule
$D$ & 1 & 2 & 3 & 5 & 10 & 20 & 50\\
\midrule
$V([D])$ & 1.98 & 1.66 & 1.17 & 0.59 & 0.21 & 0.09 & 0.03\\
$\min_{|S|=2}V(S)$ & -- & 1.66 & 1.65 & 1.53 & 1.39 & 1.40 & --\\
\bottomrule
\end{tabular}
\caption{Both classes $N(0,I_D)$, $n_0=n_1=15$, hinge loss,
$\rho_0+\rho_1=0.6$; mean over five draws. Development check.}
\label{tab:noise}
\end{table}

The detector certificate already contains what an honest radius needs.

\begin{proposition}\label{prop:honest}
Let $(g,a,\lambda)$ solve the detector program at $S$, and let $f$ be
its extension in Appendix~A of the theory note, used as the test
that declares class~$1$ when $f<0$. For any distributions
$\mathbb P_0,\mathbb P_1$ on $\mathbb R^D$ with finite first moments,
\[
\mathbb P_0(f<0)+\mathbb P_1(f\ge0)
\le\mathbb E_{\mathbb P_0}\ell(-f)+\mathbb E_{\mathbb P_1}\ell(f)
\le V(S)+\sum_{k=0}^1\lambda_k\bigl(\mathcal W_1(\widehat{\mathbb P}_k[S],
\mathbb P_k[S])-\rho_k\bigr).
\]
\end{proposition}
\begin{proof}
The first inequality holds because $\ell\ge1$ on $[0,\infty)$. The
envelope $E_k$ of Appendix~A dominates the class-$k$ loss of $f$, is
$\lambda_k$-Lipschitz for the $\ell_1$ norm on $\mathbb R^{|S|}$, and
satisfies $E_k(x_i^{(k)}[S])\le a_i^{(k)}$. Kantorovich--Rubinstein
duality gives $\mathbb E_{\mathbb P_k}E_k\le\frac1{n_k}\sum_ia_i^{(k)}
+\lambda_k\mathcal W_1(\widehat{\mathbb P}_k[S],\mathbb P_k[S])$.
Summing over $k$ gives the second inequality.
\end{proof}

So $V(S)$ bounds the true total error as soon as each projected
empirical distribution lies within $\rho_k$ of the truth. Coverage is
needed only in dimension $|S|$, and only for the subsets under
consideration. For data in $[0,1]^D$, the rates of empirical
measures~\cite{fournier2015rate,weed2019sharp}, a bounded-difference
inequality and a union bound over the at most $(eD/d)^d$ subsets
suggest a radius of order
\[
\rho_k(d,\delta)=d\,n_k^{-1/d}
+d\sqrt{\bigl(d\log(eD/d)+\log(1/\delta)\bigr)/n_k}\qquad(d\ge3),
\]
which makes the bound hold for all $|S|\le d$ at once with probability
at least $1-\delta$. Three results then seem within reach.
\begin{enumerate}
\item \emph{Validity after selection.} The selected $\widehat S$ and its
detector satisfy the error bound $V(\widehat S)$ with no correction for
the data-driven choice.
\item \emph{A budget chosen by the certificate.} The bound is uniform
over $d\le d_{\max}$, so
$\widehat d=\arg\min_d\min_{|S|=d}V_{\rho(d)}(S)$ keeps a valid
certificate. More coordinates separate the classes better but need
larger radii, and the certificate balances the two.
\item \emph{Recovery.} If $\mathbb P_0$ and $\mathbb P_1$ differ only on
a set $S^\star$, with a margin in robust total variation, then
$\widehat S=S^\star$ once $n_k$ is large compared with the radius in
dimension $|S^\star|$. The rate depends on $d$ rather than $D$, as in
the power analyses of~\cite{wang2021projected,wang2023variable}.
\end{enumerate}
A sharper route replaces coverage by the variation-based analysis of
Wasserstein robust optimization~\cite{gao2023finite}; the Lipschitz
constants $\lambda_k$ in \hyperref[prop:honest]{Proposition~\ref*{prop:honest}}
play the role of the variation used there. A further question links back to
\Secref{sec:gap}: a dense weighted metric has a large effective
dimension and would need a large honest radius, so with honest radii
the weighted-metric relaxation may lose its advantage and become
tighter. The story is that selection is not only for interpretability:
it is what makes a Wasserstein robust test valid in high dimension.

\subsection{Selection as matching interdiction}\label{sec:interdiction}
By \Secref{sec:gap}, the defender prices $d$ coordinates, and the
adversary buys the largest cross-class matching it can afford. This is
an interdiction game in which the defender raises costs and the
follower solves a matching program; arc removal is the classical
case~\cite{zenklusen2010matching}. The follower's feasible set shrinks
as more coordinates are priced, which is close to the monotone setting
of the interdiction cuts of~\cite{fischetti2019interdiction}.

The immediate step is to replace outer approximation by
\Eqref{eq:milp} for the hinge loss: it is exact, has $n_0n_1$ rows and
one price, and its relaxation is stronger than the current root.
\Tabref{tab:gap} also shows its limit. The search tree grows almost as
fast as enumeration, as \Secref{sec:gap} predicts, so the remaining
gains come from structure.
\begin{enumerate}
\item \emph{Bounds from matchings.} A matching $\pi$ stays feasible
after scaling, so for every $S$
\[
m_S(\rho)\ge\sum_{i,j}\pi_{ij}\cdot\min\Bigl\{1,\frac{\rho}{\sum_{m\in S}\kappa_m}\Bigr\},
\qquad \kappa_m=\sum_{i,j}\pi_{ij}C^m_{ij}.
\]
The bound uses the exact cost of $\pi$ under $S$; the cuts of the
theory note are weaker linear bounds of the same kind. Such bounds in
the master problem, with the pair rows of \Eqref{eq:milp} generated
lazily, give the natural decomposition for large $n$: only pairs that
are cheap in some candidate $S$ matter.
\item \emph{A statistical tolerance.} By
\hyperref[prop:honest]{Proposition~\ref*{prop:honest}}, differences in
$V$ below the slack of the certificate carry no statistical meaning.
Stopping the search at that tolerance turns exact certification into
honest certification, and may stop early.
\item \emph{General losses.} The same construction gives a compact
program with two prices and $N^2$ rows. For the hinge loss,
\hyperref[prop:matching]{Proposition~\ref*{prop:matching}} reduces it
to one price and $n_0n_1$ rows.
\end{enumerate}

\subsection{Randomized selection}\label{sec:random}
The relaxation gap measures how much mixing coordinates helps the
defender. A randomized test collects this gain while each realization
still reads only $d$ coordinates. Let the defender draw a random pair
$(S,f)$ with $|S|\le d$ and $f\in\mathcal F_S$, and let the adversary
choose $Q_k\in\mathcal P_k$ knowing the distribution of the pair but
not its realization. The risk is linear in both strategies, so
Sion's theorem~\cite{sion1958minimax}, under the usual compactness
conditions, gives
\[
\min_{(S,f)\ \text{random}}\ \max_{Q_k\in\mathcal P_k}
\mathbb E\bigl[\mathbb E_{Q_0}\ell(-f)+\mathbb E_{Q_1}\ell(f)\bigr]
=\max_{Q_k\in\mathcal P_k}\min_{|S|\le d}
\bigl[2-D_\ell(Q_0[S],Q_1[S])\bigr].
\]
The right side asks for one attack that confuses every
$d$-coordinate projection at once, so the adversary must split its
budget across coordinates. For example, take one point per class,
$x^{(0)}=(0,0)$ and $x^{(1)}=(1,1)$, the hinge loss,
$\rho_0+\rho_1=1/2$ and $d=1$. Each fixed coordinate gives $V=1$. On
the right side, an overlap $o_m$ of the projections on coordinate $m$
costs at least $o_m$, and the $\ell_1$ cost adds over coordinates, so
$o_1+o_2\le1/2$ and the value $\min_m2o_m$ is at most $1/2$: half the
fixed-selection optimum.

Randomization is known to strictly help in adversarial
classification~\cite{pinot2020randomization}; here the exchange of
minimum and maximum is exact, and the value is a divergence between
projections. Two questions come first: whether the finite reduction of
Lemma~2.1 survives in the mixed game, and how to compute the right
side, for instance by generating subsets whose pricing problem is a
selection for a fixed attack, with no inner maximization. The
guarantee holds on average over the defender's coin, which not every
application accepts.

\section{Priorities}\label{sec:priorities}
\begin{enumerate}
\item \emph{Honest radii} (\Secref{sec:honest}) offer the strongest
story at the lowest cost.
\hyperref[prop:honest]{Proposition~\ref*{prop:honest}} is already
proved; the next result is the uniform radius with a budget chosen by
the certificate, and \Tabref{tab:noise} is the motivating figure.
\item \emph{Matching interdiction} (\Secref{sec:interdiction}) starts
with a direct implementation of \Eqref{eq:milp} and a comparison on the
100-dimensional study, then adds the matching bounds and the
statistical tolerance.
\item \emph{Randomized selection} (\Secref{sec:random}) is exploratory;
the finite reduction for the mixed game comes first.
\end{enumerate}
Three lines are not pursued. Further cut engineering for outer
approximation, such as cheap least-favorable cuts, Pareto cuts or
several tangents per evaluation, reduced oracle time by
$30$--$45\%$ in development checks but cannot close a gap that belongs
to the problem. Sinkhorn ambiguity sets place the selection inside
$\exp(-c/\varepsilon)$ and lose the matching structure. Neural
generators of least favorable distributions lose the certificate.
All numbers in this note come from small synthetic development checks,
not from formal experiments.

\begingroup\small
\bibliographystyle{plain}
\bibliography{directions}
\endgroup
\end{document}
